%!TeX program = xelatex
\documentclass{SYSUReport}

% 根据个人情况修改
\headl{FPGA原理与系统设计}
\headc{课程报告}
\headr{DE2-115 RISC-V 计算机}
\lessonTitle{FPGA原理与系统设计课程报告}
\reportTitle{基于 DE2-115 的 RISC-V 软核计算机}
\stuname{张志文}
\stuid{2330919}
\inst{电子与信息工程学院}
\major{通信工程}
\date{2026年6月}

\begin{document}

% =============================================
% Part 1: 封面
% =============================================
\cover
\thispagestyle{empty}
\clearpage

% =============================================
% Part 2： 摘要
% =============================================
\begin{abstract}

本项目在 Altera DE2-115 开发板~\cite{de2115}（Cyclone IV E, EP4CE115F29C7）上构建了一台完整的 RISC-V 计算机。
以 NEORV32 软核处理器~\cite{neorv32}为核心，通过 Wishbone 总线~\cite{wishbone}连接 12 个自定义外设（SDRAM 控制器、VGA 文本/像素终端、PS/2 键盘、音频合成器、GPU 2D 加速器等），
运行 FreeRTOS 实时操作系统~\cite{freertos}，通过 UART bootloader 从 SDRAM 执行用户固件。

系统提供 20 条 shell 命令，包括：硬件加速的 Conway 生命游戏~\cite{conway1970game}、Snake 双人游戏、TWM 像素窗口管理器、AES/SHA/SM4 密码学基准测试、NTT 数论变换加速器、
13 个课程实验的硬件多路复用器、RISC-V 汇编监视器等。所有交互式程序统一 F1 帮助 / F10 退出操作逻辑。

项目历时 12 天、110 余次提交，覆盖 RTL 设计、固件开发、RTOS 移植、硬件加速器设计。

\end{abstract}
\thispagestyle{empty}
\clearpage

% =============================================
% Part 3： 目录页
% =============================================
\pagenumbering{Roman}
\setcounter{page}{1}
\tableofcontents
\clearpage

% =============================================
% Part 4： 正文内容
% =============================================
\pagenumbering{arabic}
\setcounter{page}{1}

% ============================================================
\section{项目概述}
% ============================================================

\subsection{背景与动机}

FPGA 原理课的常规做法是用 VHDL 写一个计数器或状态机，在开发板上验证时序，交一份仿真波形。DE2-115 开发板拥有 114,480 个逻辑单元、128MB SDRAM、VGA 输出、PS/2 接口、音频 CODEC、红外接收器。一次实验只用到其中几百个逻辑单元。

我想做的是把 DE2-115 变成一台通用计算机：每个外设都是系统的一个组件，而不是一次性的实验载体。

\subsection{设计目标}

\begin{itemize}
    \item 用 RISC-V 软核替代手工状态机，FPGA 运行一个系统而不是实现一个功能
    \item 统一的 shell 框架，新功能注册一个 \texttt{program\_t} 结构体即可
    \item 课程实验成为 shell 命令，在集成环境中演示
    \item 硬件加速器解决 CPU 跑不动的场景（像素填充、数论变换、生命游戏）
    \item VGA 显示 + PS/2 键盘 + 统一帮助系统
\end{itemize}

\subsection{项目规模}

\begin{table}[!htbp]
    \centering
    \begin{tabular}{l|r}
        \hline
        \textbf{指标} & \textbf{数据} \\
        \hline
        开发周期 & 12 天 (2026-05-21 \textasciitilde{} 2026-06-02) \\
        Git 提交数 & 110+ \\
        Wishbone 外设数 & 12 \\
        Shell 命令数 & 18 \\
        固件大小 & \textasciitilde156KB \\
        RTL 文件 & 20+ VHDL 源文件 \\
        C 源文件 & 30+ 源文件 \\
        FPGA 资源占用 & \textasciitilde{}30\% 逻辑单元 \\
        \hline
    \end{tabular}
    \caption{项目规模统计}
\end{table}

% ============================================================
\section{系统架构}
% ============================================================

\subsection{硬件架构}

系统顶层实体为 \texttt{de2os\_top}，核心架构如图~\ref{fig:arch}所示。

\begin{figure}[!htbp]
    \centering
    \begin{verbatim}
                  +-----------------------+
                  |     de2os_top         |
                  |  (board-level pins)   |
                  +-----------+-----------+
                              |
                  +-----------v-----------+
                  |   clk_rst_gen (PLL)   |
                  | 50MHz + 100MHz + 25MHz|
                  +-----------+-----------+
                              |
                  +-----------v-----------+
                  |  neorv32_wrapper      |
                  |  RV32IMC + Zicsr/Zbk*~\cite{riscvspec} |
                  |   IMEM 64KB / DMEM 16KB  |
                  +-----------+-----------+
                              | XBUS (Wishbone)
                  +-----------v-----------+
                  |   wb_intercon         |
                  | 1-master, 12-slave    |
                  +---+---+---+---+---+---+
                      |   |   |   |   |
           +----------+   |   |   |   +----------+
           |              |   |   |              |
    +------v------+  +----v---+   |       +------v------+
    | sdram_ctrl  |  | vga_   |   |       | conway_    |
    | 128MB@100MHz|  | text_  |   |       | engine     |
    |             |  | terminal|  |       | 64x25 grid |
    +-------------+  +--------+   |       +------------+
                                  |
    +-----------+  +----------+  |  +------------+  +-----------+
    | ps2_ctrl  |  | ntt_sdf  |<--+  | synth_    |  | gpu_2d    |
    | keyboard  |  | NTT accel|     | engine    |  | FILL rect |
    +-----------+  +----------+     | 3xOSC+FM  |  +-----------+
                                   +-----------+
    \end{verbatim}
    \caption{系统硬件架构}
    \label{fig:arch}
\end{figure}

\subsection{地址映射}

Wishbone 总线采用 1 主 12 从架构，地址分配如表~\ref{tab:addr}所示。

\begin{table}[!htbp]
    \centering
    \small
    \begin{tabular}{l|l|r|l}
        \hline
        \textbf{基地址} & \textbf{外设} & \textbf{大小} & \textbf{说明} \\
        \hline
        0x00000000 & IMEM   & 64KB  & M9K BRAM, bootloader \\
        0x80000000 & DMEM   & 16KB  & CPU 内部数据存储器 \\
        0x01000000 & SDRAM  & 128MB & 主程序执行 + FreeRTOS 堆 \\
        0x01800000 & 帧缓冲 & 600KB & VGA 像素模式 RGB565 \\
        0xF0000000 & VGA 终端 & 32KB & 80$\times$30 文本 + 像素模式 \\
        0xF0008000 & PS/2   & 4KB   & 键盘扫描码 + IRQ \\
        0xF000F000 & NTT    & 4KB   & 数论变换加速器 \\
        0xF0010000 & ExpDemo & 4KB  & 13 个课程实验多路复用 \\
        0xF0011000 & Conway & 4KB   & 64$\times$25 生命游戏引擎 \\
        0xF0012000 & Synth  & 4KB   & 音频合成 (3xOSC + FM) \\
        0xF0015000 & GPU 2D & 4KB   & 矩形填充加速器 \\
        \hline
    \end{tabular}
    \caption{Wishbone 地址映射}
    \label{tab:addr}
\end{table}

\subsection{软件架构}

固件运行在 FreeRTOS 上，采用 4 任务架构：

\begin{itemize}
    \item \texttt{t\_uart\_input}：轮询 UART 和 PS/2 键盘输入，将按键事件分发到活动程序
    \item \texttt{t\_shell}：Shell 主循环，解析命令并启动/停止程序
    \item \texttt{t\_active\_prog}：运行当前活动程序的 \texttt{update()} 回调
    \item \texttt{t\_status}：驱动 7 段数码管、LED、LCD 显示系统状态
\end{itemize}

Boot mode 0 工作流：IMEM 仅存放 \textasciitilde2KB bootloader，主程序通过 UART 上传至 SDRAM 执行。软件迭代周期从 Quartus 重编译的 \textasciitilde{}40 分钟降到增量编译的 \textasciitilde{}25 秒。

% ============================================================
\section{关键技术实现}
% ============================================================

\subsection{NEORV32 RISC-V 软核集成}

NEORV32 (v1.13.1)~\cite{neorv32} 是一款开源的 RISC-V 处理器核，支持 RV32IMC 指令集和多种扩展。
本项目通过 \texttt{neorv32\_wrapper} 进行平台适配：

\begin{itemize}
    \item 启用 ISA 扩展：IMC, Zicsr, Zicntr, Zbkb, Zbkc, Zbkx, Zknd, Zkne, Zknh, Zksed, Zksh
    \item 关闭 Zfinx（综合资源开销过大）、ICACHE（burst CDC 未完成验证）
    \item XBUS Wishbone 主接口超时设为 2048 周期（\textasciitilde{}41$\mu$s @50MHz）
    \item \texttt{std\_ulogic} $\leftrightarrow$ \texttt{std\_logic} 类型转换在 wrapper 层完成
\end{itemize}

\subsection{SDRAM 控制器}

自研 100MHz SDRAM 控制器，状态机处理初始化、刷新、读写操作。
关键参数：DRAM\_CLK 需 +1.56ns 相移（经验值），支持 Wishbone burst 传输的异步 FIFO CDC 路径（预布线，ICACHE 启用后可用）。

\subsection{VGA 文本/像素双模式终端}

\texttt{vga\_text\_terminal} 模块同时支持：

\begin{itemize}
    \item \textbf{文本模式}：80$\times$30 字符，256 字符集（CP437），每字符独立前景/背景色
    \item \textbf{像素模式}：640$\times$480@60Hz RGB565，帧缓冲位于 SDRAM，由 GPU 2D 加速器或 CPU 直接写入
\end{itemize}

像素时钟由 PLL c3 输出 25MHz（全局时钟网络），替代了先前不稳定的 toggle flip-flop 方案。

\subsection{硬件加速器}

\subsubsection{Conway 生命游戏引擎}

\texttt{conway\_engine.vhd} 实现 64$\times$25 网格的 B3/S23 规则计算，环面拓扑回绕。

\begin{itemize}
    \item 双缓冲：读/写分用 grid\_a/grid\_b，每代后翻转
    \item 逐时钟计算：每周期处理一个细胞及其 8 邻居
    \item LFSR 随机填充：16 位最大长度 LFSR 生成伪随机初始状态
    \item Wishbone 从接口：CPU 通过 5 个寄存器控制引擎
    \item ramstyle="logic"：使用逻辑单元寄存器而非 M9K，支持 32 位宽并行行读取
\end{itemize}

\subsubsection{GPU 2D 加速器}

\texttt{gpu\_2d.vhd} 通过 SDRAM burst 写实现矩形区域填充。
CPU 逐像素写 framebuffer 每帧需 \textasciitilde{}2 秒，GPU burst 写降至 \textasciitilde{}3ms，加速比约 400$\times$。

\subsubsection{NTT 数论变换加速器}

\texttt{ntt\_sdf.vhd} 实现 ML-KEM-512 所需的数论变换。
CPU 执行一轮约 10ms，硬件约 100$\mu$s，加速比约 100$\times$。

\subsection{音频合成器}

\texttt{synth\_engine.vhd} 实现双模式音频合成：

\begin{itemize}
    \item \textbf{3xOSC}：三振荡器叠加，支持锯齿波/方波/三角波/正弦波
    \item \textbf{DX7 FM}：6 操作器频率调制合成，兼容 DX7 算法
    \item I2S 输出至 WM8731 CODEC，PS/2 键盘映射 88 音符
\end{itemize}

\subsection{FreeRTOS 移植}

将 FreeRTOS~\cite{freertos} 移植到 NEORV32 裸机环境：

\begin{itemize}
    \item 定时器源：NEORV32 内置mtime（CLINT），1ms tick
    \item 堆管理：FreeRTOS heap 位于 SDRAM (0x01900000)，大小 16KB
    \item 优先级调度：4 任务固定优先级，无动态优先级调整
    \item 输入多路复用：PS/2 键盘和 UART 同时输入，通过 FreeRTOS 队列分发
\end{itemize}

% ============================================================
\section{课程实验集成}
% ============================================================

\texttt{expdemo} 模块通过 Wishbone MMIO 驱动硬件实验多路复用器 (\texttt{expdemo\_top.vhd})，
将 13 个课程实验统一为一个 shell 命令。

\begin{table}[!htbp]
    \centering
    \small
    \begin{tabular}{c|l|l}
        \hline
        \textbf{编号} & \textbf{实验名称} & \textbf{功能} \\
        \hline
        1  & LED 流水灯       & 8 种 LED 闪烁模式 \\
        2  & 按键消抖         & 硬件消抖 + 计数器 \\
        3  & 数码管动态扫描   & 8 位 7 段数码管复用 \\
        4  & 矩阵键盘扫描     & 4$\times$4 矩阵键盘检测 \\
        5  & PWM 调光         & LED 呼吸灯效果 \\
        6  & 蜂鸣器音乐       & PWM 驱动蜂鸣器播放旋律 \\
        7  & 计时器           & 精确计时 + 数码管显示 \\
        8  & 计数器           & 可重置计数器 \\
        9  & 红外解码         & NEC 协议红外遥控解码 \\
        10 & VGA 显示         & 彩条/图案生成 \\
        11 & PS/2 键盘        & 扫描码读取与显示 \\
        12 & 频率计           & 数字信号频率测量 \\
        13 & RS232 通信       & UART 收发回环 \\
        \hline
    \end{tabular}
    \caption{13 个课程实验}
    \label{tab:exp}
\end{table}

每个实验由 \texttt{expdemo\_top} 内部的独立实验适配器实现，通过 Wishbone 寄存器选择当前活跃实验。
用户通过 PS/2 键盘输入编号切换实验，VGA 终端显示实验状态。

这样课程实验不再是独立的 .sof 文件。验收时在 shell 中运行 \texttt{expdemo}，逐一演示 13 个实验。

% ============================================================
\section{软件功能}
% ============================================================

\subsection{Shell 命令}

系统提供 20 条 shell 命令，分为程序命令和工具命令两类。

\begin{table}[!htbp]
    \centering
    \small
    \begin{tabular}{l|l}
        \hline
        \textbf{命令} & \textbf{说明} \\
        \hline
        \multicolumn{2}{c}{\textit{交互式程序}} \\
        \hline
        hello    & LED 跑马灯 + 计数器 \\
        crypto   & AES/SHA/SM4 加密基准测试 \\
        ps2      & PS/2 键盘测试 \\
        snake    & Snake 双人游戏 (WASD + 方向键) \\
        conway   & Conway 生命游戏 (硬件加速 64$\times$25) \\
        info     & 系统仪表盘 \\
        riscvasm & RISC-V 寄存器/内存监视器 \\
        expdemo  & 13 个课程实验 \\
        twm      & TWM 像素窗口管理器 \\
        ntt      & NTT 加速器测试 \\
        synth    & 音频合成器 (PS/2 钢琴) \\
        \hline
        \multicolumn{2}{c}{\textit{工具命令}} \\
        \hline
        clear     & 清屏 \\
        selfcheck & 板卡自检 (启动时自动运行) \\
        stats     & 任务/堆/CPU 状态 \\
        ver       & 版本/构建信息 \\
        pxtest    & VGA 像素模式诊断 \\
        vgadump   & VGA 帧缓冲转储 \\
        vgamon    & 周期性 VGA 状态监控 \\
        \hline
    \end{tabular}
    \caption{Shell 命令列表}
    \label{tab:cmds}
\end{table}

\subsection{统一交互规范}

所有交互式程序遵循统一的操作逻辑：

\begin{itemize}
    \item \textbf{F1}：显示帮助叠加层（游戏自动暂停），再次按 F1 关闭
    \item \textbf{F10}：帮助打开时关闭帮助；帮助关闭时退出到 shell
    \item \textbf{WASD / 方向键}：移动光标/角色
    \item \textbf{Enter}：确认/运行/暂停
    \item \textbf{空格}：切换/操作
\end{itemize}

学会一个程序的操作就能操作所有程序，不需要对每个程序单独记忆。

% ============================================================
\section{开发过程}
% ============================================================

\subsection{时间线}

\begin{table}[!htbp]
    \centering
    \small
    \begin{tabular}{c|l|l}
        \hline
        \textbf{日期} & \textbf{阶段} & \textbf{关键成果} \\
        \hline
        Day 1 (05-21) & 项目初始化 & NEORV32 上板，UART Hello World \\
        Day 2 (05-22) & 引脚修正 & Pin assignment 修复 \\
        Day 3 (05-23) & 骨架搭建 & SDRAM + VGA + PS/2 + Shell + 13 实验 + 加速器 RTL \\
        Day 4 (05-24) & V2 验证 & V2 上板通过 \\
        Day 5 (05-25) & V3 架构 & V2 冻结，FreeRTOS 迁移，boot mode 0 \\
        Day 6 (05-26) & 调试验证 & VGA 调试，ExpDemo 路由 \\
        Day 7-8 (05-29/30) & 外设集成 & NTT/PONG/Conway RTL，Audio synth \\
        Day 9 (05-30) & 游戏功能 & Snake 2P，PS/2 TUI 虚拟键码 \\
        Day 10-11 (06-01) & 像素 GUI & GPU 2D，VGA PLL 修复，TWM 上板 \\
        Day 12 (06-02) & 打磨收尾 & F1/F10 统一，Conway 修复，文档完善 \\
        \hline
    \end{tabular}
    \caption{开发时间线}
\end{table}

\subsection{关键决策}

项目过程中做出的 7 个决策，每个解决一个具体瓶颈：

\begin{enumerate}
    \item \textbf{RISC-V 软核替代手工状态机}：FPGA 只用了 1\% 资源，实验之间没有关联
    \item \textbf{统一 shell 框架先建壳}：每个新功能都要从头搭框架
    \item \textbf{课程实验变成 shell 命令}：实验做完就忘，无法展示集成效果
    \item \textbf{引入 FreeRTOS + boot mode 0}：IMEM 64KB 放不下大程序，改代码不再触发 Quartus 重编译
    \item \textbf{硬件加速器针对 CPU 瓶颈}：CPU 逐像素写帧 2 秒，NTT 运算慢两个数量级
    \item \textbf{VGA 像素模式 + TWM}：文字终端和串口输出没有视觉区分度
    \item \textbf{统一操作逻辑，打磨体验}：各程序操作不一致，Bug 影响可用性
\end{enumerate}

\subsection{遇到的典型问题}

\begin{table}[!htbp]
    \centering
    \small
    \begin{tabular}{l|l|l}
        \hline
        \textbf{问题} & \textbf{根因} & \textbf{解决方案} \\
        \hline
        Conway 网格全零 & M9K BRAM 不支持并行读取 & ramstyle="logic" (LE 寄存器) \\
        SDRAM 不稳定 & DRAM\_CLK 相移不足 & +1.56ns PLL 相移 (经验值) \\
        INTC 总线挂死 & ack 信号未回传 & ack loopback \\
        I2C 总线竞争 & SDA 方向冲突 & 三态缓冲 inout \\
        VGA 文本重影 & 像素时钟抖动 & PLL c3 全局时钟网络 \\
        Snake 缓冲区溢出 & body shift 未 clamp & 循环上限 clamped to len-1 \\
        ICACHE burst 死锁 & stb 相关的 FIFO pop & 协议修正 (设计中) \\
        \hline
    \end{tabular}
    \caption{典型问题与解决方案}
\end{table}

% ============================================================
\section{工具链与构建}
% ============================================================

\subsection{开发环境}

\begin{table}[!htbp]
    \centering
    \begin{tabular}{l|l}
        \hline
        \textbf{工具} & \textbf{版本/路径} \\
        \hline
        Quartus Prime & 23.1std Lite \\
        NEORV32 & v1.13.1 (release tag) \\
        RISC-V GCC & Docker 镜像 de2extra-builder \\
        串口 & COM10, 115200 8N1 \\
        FPGA & DE2-115, EP4CE115F29C7 \\
        \hline
    \end{tabular}
    \caption{开发环境}
\end{table}

\subsection{构建流程}

提供一条龙部署脚本：

\begin{itemize}
    \item \texttt{./run/deploy\_de2shell\_rtos.sh inc}：增量编译（仅 C 代码，\textasciitilde{}25 秒）
    \item \texttt{./run/deploy\_de2shell\_rtos.sh full}：全量编译（C + VHDL + Quartus + 烧录 + 上传，\textasciitilde{}40 分钟）
    \item \texttt{./run/deploy\_de2shell\_rtos.sh fpga}：仅 RTL（bootloader + Quartus + 烧录）
\end{itemize}

Boot mode 0 把硬件和软件的迭代周期分开：改 C 代码走增量编译，不触发 Quartus。

% ============================================================
\section{AI 辅助使用说明}

本项目开发过程中使用了 Claude Code (Anthropic) 作为编程辅助工具。

\subsection{使用范围}

\begin{table}[!htbp]
    \centering
    \small
    \begin{tabular}{l|l|l}
        \hline
        \textbf{环节} & \textbf{AI 参与程度} & \textbf{说明} \\
        \hline
        架构设计 & 辅助分析 & 评估 FreeRTOS vs bare-metal、boot mode 选型 \\
        RTL 编写 & 生成初版 + 人工审查 & VHDL 模块由 AI 生成，人工审查调试 \\
        C 固件开发 & 生成初版 + 调试 & 驱动程序、shell 框架、游戏逻辑 \\
        Bug 诊断 & 辅助定位 & 分析寄存器状态、提出排查方向 \\
        文档撰写 & 辅助生成 & README、阶段文档、本报告初稿 \\
        工具链配置 & 配置辅助 & QSF 引脚分配、编译脚本、Docker \\
        \hline
    \end{tabular}
    \caption{AI 使用范围}
\end{table}

\subsection{使用原则}

\begin{enumerate}
    \item AI 生成的代码全部经过人工审查和上板验证
    \item 架构决策由人做出（引入 FreeRTOS、boot mode 0、加速器选型等）
    \item 调试以人工为主：示波器测量、上板验证、逻辑分析
    \item 对 AI 生成的每段代码，开发者都能解释其工作原理
\end{enumerate}

\subsection{AI 无法替代的工作}

\begin{itemize}
    \item 上板验证：在物理板卡上运行固件，观察 LED、VGA、键盘的实际行为
    \item 时序调试：SDRAM 相移、PLL 配置需要实际测量和经验判断
    \item 资源约束：FPGA LE 使用率、M9K 分配、时序收敛
    \item 知识内化：理解 RISC-V 指令集、Wishbone 协议、VGA 时序等底层原理
\end{itemize}

\subsection{实际效果}

AI 加速了重复性代码的编写（Wishbone 从接口、shell 命令注册、VGA 字符输出），从小时级降到分钟级。也帮助快速定位常见错误（信号未连接、类型不匹配、边界条件遗漏）。

但 AI 也引入了额外成本：生成的代码需要仔细审查。Conway 引擎的 M9K ramstyle 问题就是 AI 引入的隐蔽 bug，排查花的时间比手写还长。

\tbox{
    AI 在本项目中是"高级代码补全"和"知识检索"工具。核心设计决策、验证和知识内化都由人完成。
}

% ============================================================
\section{项目成果与展望}
% ============================================================

\subsection{成果总结}

12 天内完成的工作量：

\begin{itemize}
    \item 12 个 Wishbone 外设（含 3 个硬件加速器）
    \item FreeRTOS 实时操作系统（4 任务调度）
    \item 20 条 shell 命令
    \item VGA 文本 + 像素双模式图形系统
    \item PS/2 键盘 + UART 双路输入
    \item 13 个课程实验的硬件多路复用
\end{itemize}

\subsection{未来方向}

\begin{itemize}
    \item pForth 编译器集成：在板卡上运行 Forth 语言
    \item TWM 窗口管理器 API 开放：允许第三方程序使用像素 GUI
    \item ChromaShader 集成：实时视频色彩处理
    \item ICACHE 启用：利用预布线的 burst CDC 路径提升 SDRAM 访问性能
    \item VGA 文本重影修复：PLL 时钟路径优化
\end{itemize}

% =============================================
% Part 5： 参考文献
% =============================================

\newpage
\bibliographystyle{ieeetr}
\bibliography{reference.bib}

\end{document}
